\documentclass[10pt,a4paper]{amsart}
\usepackage[margin=19mm]{geometry}
\usepackage{amsmath,amssymb,mathtools}
\usepackage[scr=rsfs,cal=euler]{mathalfa}
\usepackage{xfrac}
\usepackage[unicode,hidelinks,bookmarks=false]{hyperref}
\newcommand{\ie}{i.e.\ }
\newcommand{\cf}{cf.\ }
\newcommand{\resp}{respectively\ }
\newcommand{\Subsection}{\subsection}
\providecommand{\nobreakdash}{\nobreak\hbox{-}}
\setlength{\emergencystretch}{2em}
\multlinegap0pt
\usepackage{mathtools}
\usepackage{amssymb}
\usepackage{xcolor}
\newtheorem{theorem}{Theorem}
\def \N {\mathbb{N}}
\def \Ns {\mathscr{N}_{s}}
\def \R {\mathbb{R}}
\renewcommand{\geq}{\geqslant}
\renewcommand{\le}{\leqslant}
\renewcommand{\leq}{\leqslant}
\title{The Neumann condition\\ for the superposition of fractional Laplacians}
\author{Serena Dipierro, Edoardo Proietti Lippi, Caterina Sportelli, Enrico Valdinoci}
\date{Source: arXiv:2402.05514v2 (CC BY 4.0)}
\begin{document}
\maketitle

\noindent\emph{Source and licence.} The Neumann condition\\ for the superposition of fractional Laplacians. Version arXiv:2402.05514v2 (CC BY 4.0), distributed by arXiv under the Creative Commons Attribution 4.0 International licence, http://creativecommons.org/licenses/by/4.0/, as recorded in the arXiv licence metadata for this exact version and shown on the abstract page https://arxiv.org/abs/2402.05514v2. Full licence notice: \texttt{licenses/arxiv-2402.05514-CC-BY-4.0.txt}.\par\smallskip

\noindent\textit{Editorial scope.} Counted unit: the article's theorem propcontinuità (source0.tex lines 295--302). The article's complete original proof of it is delivered with it (source0.tex lines 1872--2077, one \texttt{proof} environment of 206 lines, which the article places away from the statement). No mathematics is written, repaired or re-derived: the statement and the proof are the article's own lines, in the article's own notation, and no retained line is substituted or re-flowed.  No further source-local context is needed: every cross-reference of every retained line resolves inside the two delivered spans.  One declared adaptation: exactly 1 ancillary strings inside the retained spans are omitted (image-inclusion commands and, where the source repeats a label, the repeated label definitions) and no other line differs from the source; the figure environments, with their placement, captions and labels, are kept, and the package records the omitted strings in fidelity receipts bound to the affected bodies.  No retained line contains a citation, so the wrapper carries no bibliography and no bibliography entry.  The retained lines contain the article's own diagram code, which the wrapper typesets with the standard tikz packages; no image file is delivered and the package declares no asset dependency.  Editorial formatting only, in addition: an AMS \texttt{amsart} preamble that restates the article's own declarations and macros used by the retained lines, namely the environments \texttt{theorem} on separate counters, together with the standard packages those lines need.  The retained environments print in this document as Theorem 1 in order of appearance, whereas the article prints the same items with the numbering of its own full document.  The complete native source of the article as distributed in the arXiv e-print for this exact version is bundled unchanged as \texttt{sources/154222/source0.tex}.
\begin{theorem}\label{propcontinuità}
Let $\Omega\subset \R^N$ be a domain with $C^1$ boundary. Let 
$u$ be continuous in $\overline{\Omega}$, with 
\begin{equation}\label{NNN}
\int_{(0,1)}\Ns u\,d\mu(s)=0 \mbox{ in } \R^N\setminus \overline{\Omega}.
\end{equation}
Then, $u$ is continuous in the whole of $\R^N$.
\end{theorem}

\begin{proof}[Proof of Theorem~\ref{propcontinuità}]
First, we fix $x_0\in \R^N\setminus \overline{\Omega}$ and we 
prove that~$u$ is continuous at~$x_0$. Since~$\R^N\setminus \overline{\Omega}$ is an open set, there exists~$\rho>0$ such that~$|x_0-y|\geq \rho$ for any~$y\in \Omega$.
So, if~$x\in B_{\rho/2}(x_0)$, we have that
\[
|x-y|\geq |x_0-y|-|x_0-x|\geq \frac{\rho}{2}\quad\mbox{ for any } y\in\Omega.
\]
As a concequence, if $\rho$ is small enough,
for any $x\in B_{\rho/2}(x_0)$ and $y\in \Omega$
we have that 
\[
\frac{|u(y)|+1}{|x-y|^{N+2s}}
\leq \frac{2^{N+2s}}{\rho^{N+2s}}\Big(\|u\|_{L^\infty(\overline{\Omega})}+1\Big)
\leq \frac{2^{N+2}}{\rho^{N+2}}\Big(\|u\|_{L^\infty(\overline{\Omega})}+1\Big).
\]
{F}rom this, we deduce that 
\[
\begin{split}
\int_{(0,1)}c_{N,s}\int_\Omega \frac{|u(y)|+1}{|x-y|^{N+2s}}\,dy\,d\mu(s)
&\leq \frac{2^{N+2}}{\rho^{N+2}}|\Omega|\Big(\|u\|_{L^\infty(\overline{\Omega})}+1\Big)\int_{(0,1)}c_{N,s}\,d\mu(s)\\
&\le\frac{2^{N+2}}{\rho^{N+2}}|\Omega|\Big(\|u\|_{L^\infty(\overline{\Omega})}+1\Big) \sup_{s\in (0, 1)} c_{N, s} \int_{(0, 1)} d\mu(s)\\
&<+\infty.
\end{split}
\]
Thus, by the Neumann condition in~\eqref{NNN} and the dominated convergence theorem we obtain that
\[
\begin{split}
\lim_{x\to x_0}u(x) &=
\lim_{x\to x_0} \dfrac{\displaystyle \int_{(0,1)}c_{N,s}\int_\Omega \dfrac{u(y)}{|x-y|^{N+2s}}\,dy\,d\mu(s)}{\displaystyle\int_{(0,1)}c_{N,s}\int_\Omega \dfrac{1}{|x-y|^{N+2s}}\,dy\,d\mu(s)}\\
&=\dfrac{\displaystyle \int_{(0,1)}c_{N,s}\int_\Omega \dfrac{u(y)}{|x_0-y|^{N+2s}}\,dy\,d\mu(s)}{\displaystyle\int_{(0,1)}c_{N,s}\int_\Omega \dfrac{1}{|x_0-y|^{N+2s}}\,dy\,d\mu(s)}\\
&=u(x_0).
\end{split}
\]
This proves that $u$ is continuous at any point of 
$\R^N\setminus \overline{\Omega}$.

Now we show that $u$ is continuous at any point 
$p\in \partial\Omega$.
To do so, we need a more accurate argument, since numerators and denominators
may create singularities which have to be attentively estimated to detect cancellations.
We take a sequence $\{p_k\}_k$
converging to $p$ as $k$ goes to infinity.
Let $q_k$ be the projection of $p_k$ onto $\overline{\Omega}$.
Since $p\in \overline{\Omega}$, from the minimizing property of the 
projection we have that 
\[
|p_k-q_k|=\inf_{\xi\in \overline{\Omega}}|p_k-\xi|
\leq |p_k-p|,
\]
and so
\[
\lim_{k\to +\infty}|q_k-p|
\leq \lim_{k\to +\infty}\big(|q_k-p_k|+|p_k-p|\big)
\leq \lim_{k\to +\infty} 2|p_k-p|=0.
\]
Thus, since by assumption $u$ is continuous in $\overline{\Omega}$,
we have
\begin{equation}\label{pkconvergenza1}
\lim_{k\to +\infty}u(q_k)=u(p).
\end{equation}
Now we claim that
\begin{equation}\label{pkconvergenza2}
\lim_{k\to +\infty}\big(u(q_k)-u(p_k)\big)=0.
\end{equation}
To prove this, it is enough to consider the points $p_k$ that belong
to $\R^N \setminus\overline{\Omega}$. Indeed, if 
$p_k\in \overline{\Omega}$, we have $p_k=q_k$ and so 
\eqref{pkconvergenza2} is trivially satisfied.

We define $\nu_k:=(p_k-q_k)/|p_k-q_k|$, so that $\nu_k$ is the 
exterior normal of $\Omega$ at $q_k\in \partial\Omega$.
We consider a rigid motion~$\mathcal{R}_k$ such that~$\mathcal{R}_k q_k=0$ and~$\mathcal{R}_k \nu_k=e_N=(0,\cdots,0,1)$ (see Figure~\ref{RK}).  Let~$h_k:=|p_k-q_k|$. We notice that
\begin{equation}\label{dovevapk}
h_k^{-1}\mathcal{R}_k p_k=
h_k^{-1}\mathcal{R}_k (p_k-q_k)=\mathcal{R}_k \nu_k=e_N.
\end{equation}
Then, the domain
\[
\Omega_k:=h_k^{-1}\mathcal{R}_k \Omega
\]
has a vertical exterior normal at 0 and approaches the half space
$\Pi:=\{x_N<0\}$ as $k\to \infty$.

%\vspace{-5em}
\begin{figure}[h]
\begin{center}

\end{center}
\caption{The rigid motion $\mathcal R_k$ of $\Omega$.
%%%% when the domain $\Omega_k$ is Australia :)
}
\label{RK}
\end{figure}

Now, we use the homogeneous Neumann condition at $p_k$ to obtain
\[
\begin{split}
u(p_k)-u(q_k)&=\dfrac{\displaystyle\int_{(0,1)}c_{N,s}\int_{\Omega}\dfrac{u(y)}{|p_k-y|^{N+2s}}\,dy\,d\mu(s)}{\displaystyle\int_{(0,1)}c_{N,s}\int_{\Omega}\dfrac{1}{|p_k-y|^{N+2s}}\,dy\,d\mu(s)}-u(q_k)\\
&=\dfrac{\displaystyle\int_{(0,1)}c_{N,s}\int_{\Omega}\dfrac{u(y)-u(q_k)}{|p_k-y|^{N+2s}}\,dy\,d\mu(s)}{\displaystyle\int_{(0,1)}c_{N,s}\int_{\Omega}\dfrac{1}{|p_k-y|^{N+2s}}\,dy\,d\mu(s)}\\
&=I_1+I_2,
\end{split}
\]
with 
\[
I_1:=\dfrac{\displaystyle\int_{(0,1)}c_{N,s}\int_{\Omega\cap B_{\sqrt{h_k}}(q_k)}\dfrac{u(y)-u(q_k)}{|p_k-y|^{N+2s}}\,dy\,d\mu(s)}{\displaystyle\int_{(0,1)}c_{N,s}\int_{\Omega}\dfrac{1}{|p_k-y|^{N+2s}}\,dy\,d\mu(s)}
\]
and
\begin{equation}\label{I2}
I_2:=\dfrac{\displaystyle\int_{(0,1)}c_{N,s}\int_{\Omega\setminus B_{\sqrt{h_k}}(q_k)}\dfrac{u(y)-u(q_k)}{|p_k-y|^{N+2s}}\,dy\,d\mu(s)}{\displaystyle\int_{(0,1)}c_{N,s}\int_{\Omega}\dfrac{1}{|p_k-y|^{N+2s}}\,dy\,d\mu(s)}.
\end{equation}
We observe that from the uniform continuity of $u$ in $\Omega$, we
have
\[
\lim_{k\to +\infty}\sup_{\Omega\cap B_{\sqrt{h_k}}(q_k)}
|u(y)-u(q_k)|=0.
\]
Therefore,
\begin{equation}\label{I1to0}
\lim_{k\to +\infty} |I_1|\le\lim_{k\to +\infty} \, \sup_{\Omega\cap B_{\sqrt{h_k}}(q_k)} |u(y)-u(q_k)|= 0.
\end{equation}
Moreover, using the change of variable 
$\eta:=h_k^{-1}\mathcal{R}_k y$ and recalling~\eqref{dovevapk},
we obtain
\[
\begin{aligned}
|I_2|&\leq \ \dfrac{\displaystyle\int_{(0,1)}c_{N,s}\int_{\Omega\setminus B_{\sqrt{h_k}}(q_k)}\dfrac{|u(y)-u(q_k)|}{|p_k-y|^{N+2s}}\,dy\,d\mu(s)}{\displaystyle\int_{(0,1)}c_{N,s}\int_{\Omega}\dfrac{1}{|p_k-y|^{N+2s}}\,dy\,d\mu(s)}\\
&\leq 2\|u\|_{L^\infty(\overline{\Omega})} \ \dfrac{\displaystyle\int_{(0,1)}c_{N,s}\int_{\Omega\setminus B_{\sqrt{h_k}}(q_k)}\dfrac{1}{|p_k-y|^{N+2s}}\,dy\,d\mu(s)}{\displaystyle\int_{(0,1)}c_{N,s}\int_{\Omega}\dfrac{1}{|p_k-y|^{N+2s}}\,dy\,d\mu(s)} \\
&=2\|u\|_{L^\infty(\overline{\Omega})} \ \dfrac{\displaystyle\int_{(0,1)}c_{N,s}\int_{\Omega_k\setminus B_{1/{\sqrt{h_k}}}}\dfrac{h_k^N}{|h_k\mathcal{R}_k^{-1}e_N-h_k\mathcal{R}_k^{-1}\eta|^{N+2s}}\,d\eta\,d\mu(s)}{\displaystyle\int_{(0,1)}c_{N,s}\int_{\Omega_k}\dfrac{h_k^N}{|h_k\mathcal{R}_k^{-1}e_N-h_k\mathcal{R}_k^{-1}\eta|^{N+2s}}\,d\eta\,d\mu(s)} \\
&=2\|u\|_{L^\infty(\overline{\Omega})} \ \dfrac{\displaystyle\int_{(0,1)}c_{N,s}\int_{\Omega_k\setminus B_{1/{\sqrt{h_k}}}}\dfrac{h_k^{-2s}}{|e_N-\eta|^{N+2s}}\,d\eta\,d\mu(s)}{\displaystyle\int_{(0,1)}c_{N,s}\int_{\Omega_k}\dfrac{h_k^{-2s}}{|e_N-\eta|^{N+2s}}\,d\eta\,d\mu(s)}.
\end{aligned}
\]
Now, we observe that, if 
$\eta \in \Omega_k\setminus B_{1/{\sqrt{h_k}}}$, we have
\[
\begin{aligned}&
|e_N-\eta|^{N+2s}=|e_N-\eta|^{N+s}|e_N-\eta|^s
\geq |e_N-\eta|^{N+s} (|\eta|-1)^s \\
&\qquad\qquad=|e_N-\eta|^{N+s} (h_k^{-1/2}-1)^s
\geq |e_N-\eta|^{N+s} h_k^{-s/4},
\end{aligned}
\]
for $k$ large enough, and consequently
\begin{equation}\label{I2minoreuguale}
|I_2|\leq 2\|u\|_{L^\infty(\overline{\Omega})} \ \dfrac{\displaystyle\int_{(0,1)}c_{N,s}\int_{\Omega_k}\dfrac{h_k^{-\frac{7}{4}s}}{|e_N-\eta|^{N+2s}}\,d\eta\,d\mu(s)}{\displaystyle\int_{(0,1)}c_{N,s}\int_{\Omega_k}\dfrac{h_k^{-2s}}{|e_N-\eta|^{N+2s}}\,d\eta\,d\mu(s)}.
\end{equation}
Then, setting
\begin{equation}\label{I2mMM}
m:=\min\left\{n\in \N:\, \mbox{supp}(\mu) \cap\Bigg(\left(\frac{7}{8}\right)^n,\left(\frac{7}{8}\right)^{n-1}\Bigg]\neq \emptyset\right\},
\end{equation}
we can multiply and divide~\eqref{I2minoreuguale} by $h_k^{2(7/8)^m}$ to obtain
\begin{equation}\label{I2minoreuguale7/8}
|I_2|\leq 2\|u\|_{L^\infty(\overline{\Omega})}
\dfrac{\displaystyle\int_{(0,(7/8)^{m-1}]}c_{N,s}\int_{\Omega_k}\dfrac{h_k^{2\left(\frac{7}{8}\right)^m-\frac{7}{4}s}}{|e_N-\eta|^{N+2s}}\,d\eta\,d\mu(s)}{\displaystyle\int_{(0,(7/8)^{m-1}]}c_{N,s}\int_{\Omega_k}\dfrac{h_k^{2\left(\frac{7}{8}\right)^m-2s}}{|e_N-\eta|^{N+2s}}\,d\eta\,d\mu(s)}.
\end{equation}
Recalling that $\Omega_k$ converges to the half space $\Pi$, if~$k$ is large enough, then we can suppose that~$\Omega_k \subset \R^N\setminus B_{1/2}(e_N)$. In this way, 
\[
\begin{aligned}
c_{N,s}&\int_{\Omega_k}\dfrac{1}{|e_N-\eta|^{N+2s}}\,d\eta\,d\mu(s)
\leq c_{N,s}\int_{\R^N\setminus B_{1/2}(e_N)}\dfrac{1}{|e_N-\eta|^{N+2s}}\,d\eta\,d\mu(s) \\
&\quad=c_{N,s} \,\omega_{N-1}\int_{1/2}^{+\infty}\rho^{-1-2s}\,d\rho
=\frac{c_{N,s}}{s}\,2^{2s-1}\omega_{N-1},
\end{aligned}
\]
which is uniformly bounded for $s\in (0,1)$. 

Moreover, in view of~\eqref{I2mMM} we can focus on the case $s<(7/8)^{m-1}$
(otherwise we exit the support of~$\mu$).
As a result, we see that $$h_k^{2\left(\frac{7}{8}\right)^m-\frac{7}{4}s}\to 0
\quad{\mbox{ as $k\to +\infty$}}$$ and therefore, by the
dominated convergence theorem,
\begin{equation}\label{I_2numeratore}
\lim_{k\to+\infty}
\displaystyle\int_{(0,(7/8)^{m-1}]}c_{N,s}\int_{\Omega_k}\dfrac{h_k^{2\left(\frac{7}{8}\right)^m-\frac{7}{4}s}}{|e_N-\eta|^{N+2s}}\,d\eta\,d\mu(s)=0.
\end{equation}

%%Furthermore,\[
%%\begin{aligned}
%%\int_{(0,(7/8)^{m-1}]}&c_{N,s}\int_{\Omega_k}\dfrac{h_k^{2\left(\frac{7}{8}\right)^m-2s}}{|e_N-\eta|^{N+2s}}\,d\eta\,d\mu(s)\\ 
%%&= 
%%\int_{(0,(7/8)^{m}]}c_{N,s}\int_{\Omega_k}\dfrac{h_k^{2\left(\frac{7}{8}\right)^m-2s}}{|e_N-\eta|^{N+2s}}\,d\eta\,d\mu(s)+
%%\int_{((7/8)^{m},(7/8)^{m-1}]}c_{N,s}\int_{\Omega_k}\dfrac{h_k^{2\left(\frac{7}{8}\right)^m-2s}}{|e_N-\eta|^{N+2s}}\,d\eta\,d\mu(s).
%%\end{aligned}
%%\]
We observe that, for any $s\in ((7/8)^{m},(7/8)^{m-1}]$, 
\[
\lim_{k\to +\infty} c_{N,s}\int_{\Omega_k}\dfrac{h_k^{2\left(\frac{7}{8}\right)^m-2s}}{|e_N-\eta|^{N+2s}}\,d\eta=+\infty.
\]
Hence, by Fatou's lemma and~\eqref{I2mMM},
\begin{equation}\label{I2denominatore}
\liminf_{k\to+\infty}
\int_{((7/8)^{m},(7/8)^{m-1}]}c_{N,s}\int_{\Omega_k}\dfrac{h_k^{2\left(\frac{7}{8}\right)^m-2s}}{|e_N-\eta|^{N+2s}}\,d\eta\,d\mu(s)=+\infty.
\end{equation}
Using~\eqref{I_2numeratore} and~\eqref{I2denominatore} in
\eqref{I2minoreuguale7/8}, we obtain that
\begin{equation}\label{LIMI2}
\lim_{k\to+\infty} |I_2|=0.
\end{equation}
This and~\eqref{I1to0} imply~\eqref{pkconvergenza2}, as desired.

{F}rom~\eqref{pkconvergenza1} and~\eqref{pkconvergenza2} we
conclude that
\[
\lim_{k\to+\infty}u(p_k)=u(p),
\]
that is $u$ is continuous at $p$, as desired.
\end{proof}

\end{document}
